\chapter{Theoretical Foundation}
\label{chap:theory}

\section{Financial Data and Preprocessing}

\subsection{Transaction, Balance, and Cash-Flow Model}

A valid transaction has, at minimum, a positive amount $A>0$, a
transaction type, a date, a description, a category, and a wallet within
the user's scope. PERFIN distinguishes \texttt{income},
\texttt{expense}, \texttt{transfer}, \texttt{investment\_inflow},
\texttt{investment\_outflow}, and \texttt{investment\_pnl} because each
type has a different effect on balances and reports. For a wallet $w$,
the balance change caused by an income or expense transaction is

\begin{equation}
  \Delta B_w=
  \left\{
  \begin{array}{ll}
    +A, & \mbox{if income},\\
    -A, & \mbox{if expense}.
  \end{array}
  \right.
\end{equation}

Within a period $P$, the total income $I_P$, total expense $E_P$, and
net cash flow $N_P$ are computed from the transactions that have not
been soft-deleted:

\begin{equation}
  I_P=\sum_{t\in P,\,type(t)=income}A_t,\qquad
  E_P=\sum_{t\in P,\,type(t)=expense}A_t,\qquad
  N_P=I_P-E_P.
\end{equation}

Transferring $A$ from a source wallet $s$ to a destination wallet $d$
produces $\Delta B_s=-A$ and $\Delta B_d=+A$, so that
$\Delta(B_s+B_d)=0$. This invariant keeps transfers from being counted
as income or expense. The prototype allows a negative wallet balance to
reflect cash or an overdrawn account; the business condition is a
positive amount, two distinct wallets, and common ownership, not a
``sufficient balance'' check. The debit, the credit, and the history
record must live inside the same \texttt{BEGIN/COMMIT} transaction; if
one step fails, \texttt{ROLLBACK} restores the entire state
\cite{postgresql2026}.

\subsection{Integrity and Data Life Cycle}

PERFIN applies four layers of protection:

\begin{enumerate}
  \item \textbf{Domain constraints:} transaction amounts, limits, and
  contribution amounts must be valid; dates must not exceed the business
  domain.
  \item \textbf{Referential integrity:} transactions reference the user,
  the category, and the wallet through foreign keys.
  \item \textbf{Business integrity:} no transfer within the same wallet;
  no transaction written with a missing required field; a soft delete
  never destroys the ability to recover.
  \item \textbf{Atomicity:} every multi-table change or balance change
  with an accompanying history entry is committed as a whole or not at
  all.
\end{enumerate}

Indexes are placed on frequently filtered columns such as user, date,
category, wallet, and status. Money is stored as \texttt{DECIMAL(15,2)};
when carried into JavaScript it must be converted in a controlled way to
avoid string concatenation or unintended rounding. The cache is
invalidated after commit. A cache error after commit must not be
surfaced as a business error, because the client may retry and cause a
duplicated effect.

\subsection{Time-Axis Normalization and Observation Windows}

Analysis by day, week, or month must preserve the correct calendar
axis. Let $x_j$ be the total expense of the $j$-th calendar period in a
window $W$. If that period has no transactions then $x_j=0$; the period
must not be dropped from the series. Zero-filling prevents two
distortions: compressing the time gap inflates the slope, and taking
only days with spending inflates the burn rate.

The observation window must always accompany the result. PERFIN
currently uses 14 days for runway, 30 days for anomaly, 90 days for
recurring charges, 6 months for trend, and 12 weeks for correlation by
default. These are starting parameters for the prototype, not thresholds
proven optimal across every user.

\subsection{Persistent Data and Transient State}

PostgreSQL stores business data that must be auditable. Redis stores
short-lived state: pending transactions, the field currently being
clarified, ambiguous candidate lists, caches, and rate-limit counters.
A TTL expires conversational state through the \texttt{EXPIRE} mechanism
\cite{redisexpire}. The in-process memory fallback serves development
only; its data is lost on restart and it is not a production
configuration.

\section{Input Normalization and Classification}

\subsection{Normalizing Amounts, Dates, and Transaction Structure}

Natural input is turned into a typed draft before validation. The
suffixes \texttt{k}, ``ngh\`in'', and ``ng\`an'' multiply the value by
$10^3$; \texttt{tr} and ``tri\d{e}u'' multiply by $10^6$. Relative dates
are resolved against the local date of the request. A sentence with
several amounts must produce an array of drafts, in which each element
has its own \texttt{amount}, \texttt{type}, \texttt{description},
\texttt{transaction\_date}, \texttt{category}, and \texttt{wallet}. The
parser or the LLM only creates the draft; the service then checks that
the number is finite and that the date, category, wallet, and ownership
are valid before building the preview.

\subsection{Text Similarity and Safe Category Selection}

Strings are stripped of diacritics, lower-cased, cleared of
non-alphanumeric characters, and collapsed on whitespace. For two
normalized strings $a,b$, the edit score uses the Levenshtein distance
$D_{lev}$ \cite{levenshtein1966}:

\begin{equation}
  s_{edit}=1-\frac{D_{lev}(a,b)}{\max(|a|,|b|)}.
\end{equation}

Let $T_a,T_b$ be the two token sets; the Dice coefficient
\cite{dice1945} is

\begin{equation}
  s_{dice}=\frac{2|T_a\cap T_b|}{|T_a|+|T_b|}.
\end{equation}

If the smaller token set has at least two elements and is fully
contained in the larger one, the containment score $s_{contain}$ lies in
$[0.82, 0.94]$ according to the length ratio. The combined score used by
the code is

\begin{equation}
  s=\max(s_{edit},\;0.92\,s_{dice},\;s_{contain}).
\end{equation}

The selection order is exact category name, exact alias, and only then
fuzzy. Long input uses the threshold $\tau=0.82$; input of at most four
characters uses $\tau=0.90$. The best candidate must also beat the
second by at least a margin $m=0.08$. If $s_1<\tau$ or $s_1-s_2<m$, the
system chooses ``Other'' or asks for clarification instead of assigning
an ambiguous result. This formula serves FR-04 and does not replace the
user's confirmation.

\subsection{Learning from Feedback}

When a user corrects a category, the system stores the pair of the
original result and the correct result together with its context. An
exact/fuzzy correction is checked before the alias and the LLM on the
next occurrence; a conflicting correction lowers the confidence or
requests confirmation. This is retrieval combined with an adaptive rule,
\textbf{not fine-tuning} of the model. The effect of the feedback must
be evaluated by accuracy or macro-F1 before and after on the same replay
set.

\subsection{OCR, Speech-to-Text, and Provider Adapters}

OCR consists of image preprocessing, text-region detection,
recognition, and post-processing \cite{smith2007}. STT converts audio
into a transcript \cite{radford2022}. In PERFIN, a provider returns only
raw text or a transcript together with the provider name and an error
status; the result then continues through the same
draft--validation--preview pipeline as text.

A receipt image may create one aggregate transaction or several
line-item drafts, but the current version \textbf{has no algorithm that
automatically reconciles} the line-item total with the receipt total,
taxes, and discounts. The user must review the preview; a two-image
smoke test is not an OCR-accuracy measurement. The transcript must also
be shown before confirmation so that a recognition error does not
silently become an amount. The adapter allows a local provider to be
swapped for a cloud provider without changing the transaction service.
When a provider fails, the system returns an error or allows manual
entry; it does not fabricate sample data.

\section{Analytical and Planning Algorithms}

\begin{table}[htbp]
  \centering
  \caption{Deterministic algorithms and the features that use them}
  \label{tab:algorithms}
  \small
  \begin{tabularx}{\textwidth}{|p{3.0cm}|X|X|}
    \hline
    \textbf{Algorithm} & \textbf{Input and output} &
    \textbf{Feature that uses it} \\ \hline
    Linear regression & Monthly expense series with zero-fill; slope,
    $R^2$, forecast & FR-07: trend \\ \hline
    Z-score and IQR & Daily or per-item expense series; flag, score,
    method & FR-07: anomaly \\ \hline
    Cash-flow runway & Total balance and 14 calendar days; burn rate,
    days left, depletion date & FR-07/FR-11: cash-flow alert \\ \hline
    Recurring charges & Description, amount, date; cluster, cadence,
    monthly estimate & FR-07/FR-11: subscription scan \\ \hline
    Pearson & Two weekly series with zero-fill; positive correlation &
    FR-07: category relationship \\ \hline
    Budget recommender & Monthly spend, income, limit; recommendation
    and overspend forecast & FR-09 \\ \hline
    Goal planner & Target, contribution or payment, deadline, interest;
    time, shortfall, what-if & FR-10 \\ \hline
  \end{tabularx}
\end{table}

\subsection{Linear Regression and Trend Forecasting}

Let $y_i$ be the total expense of period $i$, $x_i=i$ for
$i=0,\ldots,n-1$, and $\bar{x}$ and $\bar{y}$ the respective means. The
least-squares line $\hat y_i=a+bx_i$ has \cite{montgomery2012}

\begin{equation}
  b=\frac{\sum_{i=0}^{n-1}(x_i-\bar{x})(y_i-\bar{y})}
          {\sum_{i=0}^{n-1}(x_i-\bar{x})^2},\qquad
  a=\bar{y}-b\bar{x}.
\end{equation}

The goodness of fit and the next-period forecast are

\begin{equation}
  R^2=1-\frac{\sum_i(y_i-\hat y_i)^2}
                 {\sum_i(y_i-\bar y)^2},\qquad
  \hat y_n=\max(0,a+bn).
\end{equation}

When a constant series makes the denominator of $R^2$ equal to zero, the
code returns $R^2=0$ and does not publish a trend. The service adds a
rising trend to the facts only when there are at least three months,
$b>0$, $R^2\geq0.5$, and the average percentage change between
consecutive pairs with a positive predecessor is at least 10\%.
Zero-filling must be done before the regression. The result describes a
short-term linear trend; it does not model seasonality or future events.

\subsection{Anomaly Detection with Z-score and IQR}

For $n\geq4$ expense values $x_i$, the mean, the sample standard
deviation, and the z-score are

\begin{equation}
  \mu=\frac{1}{n}\sum_{i=1}^{n}x_i,\qquad
  s=\sqrt{\frac{\sum_{i=1}^{n}(x_i-\mu)^2}{n-1}},\qquad
  z_i=\frac{x_i-\mu}{s}.
\end{equation}

Quantiles are linearly interpolated over the sorted series. Let
$IQR=Q_3-Q_1$ and the upper bound $U=Q_3+kIQR$. Following the
exploratory-data-analysis convention of Tukey \cite{tukey1977}, PERFIN
flags the high-spend side when

\begin{equation}
  z_i\geq2.5\qquad\hbox{or}\qquad x_i>Q_3+1.5\,IQR.
\end{equation}

If $s=0$ then $z_i=0$; if $IQR=0$ the IQR branch raises no flag. The
output contains the value, the date or label, $z_i$, the ratio to the
mean, and the method \texttt{z}, \texttt{iqr}, or \texttt{z+iqr}. The
result only asks the user to review; it does not prove fraud, and the
thresholds must be calibrated on representative data.

\subsection{Estimating Cash-Flow Runway}

Let $B=\sum_w B_w$ be the current total balance, $W=14$ the number of
calendar days observed, and $e_j\geq0$ the spending on day $j$; a day
with no spending has $e_j=0$. The burn rate and the days remaining are

\begin{equation}
  \bar e_W=\frac{1}{W}\sum_{j=1}^{W}e_j,\qquad
  d=\left\lfloor\frac{B}{\bar e_W}\right\rfloor.
\end{equation}

If $B\leq0$, \texttt{daysLeft=0}; if $\bar e_W=0$, the system does not
forecast a depletion date (\texttt{null}). Otherwise the depletion date
is the current date plus $d$. When a payday is known, the system finds
its next occurrence in the calendar and compares it with the depletion
date. For example, with $B=2.8$ million VND and a 14-day total spend of
$1.4$ million VND, so that $\bar e=100{,}000$ VND/day, we get $d=28$
days. This is an extrapolation that keeps the recent spending rate
constant; it does not yet account for upcoming income or unusual
spending.

\subsection{Mining Recurring Charges}

Only \texttt{expense} transactions not exceeding 500,000 VND are
considered in the current configuration. After the description is
normalized, transactions with the same key are grouped into a cluster of
$m$ amounts $a_i$ and dates $t_i$. The mean amount, the relative
deviation, and the mean cadence are

\begin{equation}
  \bar a=\frac{1}{m}\sum_{i=1}^{m}a_i,\qquad
  \delta_i=\frac{|a_i-\bar a|}{\bar a},\qquad
  \bar\Delta=\frac{1}{m-1}\sum_{i=2}^{m}|t_i-t_{i-1}|.
\end{equation}

A cluster is stable when every $\delta_i\leq0.15$. A cluster is a
recurring candidate when it is stable and satisfies one of two
conditions: $20\leq\bar\Delta\leq40$ days, or at least three
occurrences. With $m=2$, the cadence must still fall within the window.
The monthly estimate is currently $\bar a$ under the assumption of one
charge per month; the total subscription amount is the sum of the
cluster estimates. Exact grouping limits false positives but may miss
variant descriptions. The result is only a ``recurring signal''; it does
not create a recurring bill or cancel a service on its own.

\subsection{Cross-Category Spending Correlation}

For two weekly expense series $X=(x_1,\ldots,x_n)$ and
$Y=(y_1,\ldots,y_n)$ on the same weekly axis, a missing week for one
category is filled with 0. The Pearson coefficient \cite{pearson1895} is

\begin{equation}
  r=\frac{\sum_i(x_i-\bar{x})(y_i-\bar{y})}
  {\sqrt{\sum_i(x_i-\bar{x})^2\sum_i(y_i-\bar{y})^2}}.
\end{equation}

The function returns 0 if $n<3$ or if one series has zero variance. The
current version considers only categories with at least four weeks of
observation and publishes only pairs with a \textbf{positive}
correlation $r\geq0.6$; negative correlation is not yet published. The
system selects the pair with the largest $r$ and must interpret it as
``co-movement in the observed data,'' not a causal relationship.

\subsection{Budget Recommendation and Forecast}

Let $s_{c,j}$ be the spending of category $c$ in month $j$ over $m$
calendar months; a month with no activity has $s_{c,j}=0$. The mean and
the base level with a buffer $\beta=5\%$ are

\begin{equation}
  \bar s_c=\frac{1}{m}\sum_{j=1}^{m}s_{c,j},\qquad
  b_c=(1+\beta)\bar s_c.
\end{equation}

With monthly income $Y$, the default framework allocates at most
$C_{need}=0.50Y$ to needs and $C_{want}=0.30Y$ to wants, and targets
$0.20Y$ for savings, following the widely used 50/30/20 guideline
\cite{warren2005}. Three strategies produce a raw limit $L_c$:

\begin{equation}
  L_c=
  \left\{
  \begin{array}{ll}
    b_c,
      & \mbox{category average},\\
    C_g\displaystyle\frac{b_c}{\sum_{k\in g}b_k},
      & \mbox{50/30/20},\\
    b_c\displaystyle\min\left(1,\frac{C_g}{\sum_{k\in g}b_k}\right),
      & \mbox{hybrid}.
  \end{array}
  \right.
\end{equation}

Here $g$ is the needs or wants group. The result is rounded to a step of
10,000 VND; a positive value smaller than one step still becomes 10,000
VND. The confidence depends on the number of active months for the
category: high with at least 6, medium with 3--5, low below 3; every
recommendation also carries a warning if the observed history is under
three months.

To forecast overspending within the month, with $S$ the amount spent
after $d_e$ days and $D$ days in the month:

\begin{equation}
  q=\frac{S}{d_e},\qquad \widehat S_D=qD,\qquad
  d_{over}=\left\lceil\frac{L-S}{q}\right\rceil.
\end{equation}

The system attaches \texttt{likely\_to\_exceed} when $\widehat S_D>L$.
The formula assumes a constant spending rate, so it is only an early
warning; any recommendation or limit change always requires the user's
confirmation.

\subsection{Goal and Debt-Payoff Planning}

For a saving or purchase goal, let $T>0$ be the target, $C\geq0$ the
current amount, $R=\max(0,T-C)$ the remaining shortfall, and $M\geq0$
the monthly contribution. If $M>0$:

\begin{equation}
  n=\left\lceil\frac{R}{M}\right\rceil.
\end{equation}

If the deadline leaves $n_D>0$ months, the required contribution is
$M_D=\lceil R/n_D\rceil$, the monthly gap is $G=\max(0,M_D-M)$, and the
shortfall at the deadline is $H=\max(0,R-Mn_D)$. When $M=0$, the
completion time is undefined; when $R=0$, the goal is complete
immediately. A what-if only reruns the same function with $M'=M+E$,
where $E$ is the freed spending, and does not modify the original plan.

For debt, let $L$ be the current balance, $i_a$ the annual interest rate
in percent, $r=i_a/(100\times12)$ the monthly rate, and $n$ the number
of months. The level payment that finishes exactly on time is the
standard amortization formula \cite{brealey2020}

\begin{equation}
  P=
  \left\{
  \begin{array}{ll}
    L/n, & r=0,\\[2mm]
    \displaystyle\frac{Lr}{1-(1+r)^{-n}}, & r>0.
  \end{array}
  \right.
\end{equation}

The simulation uses the recurrence
$L_{k+1}=\max(0,L_k(1+r)-P)$ and adds $L_kr$ to the total interest. If
$P\leq Lr$ in the first month, the balance does not decrease and the
system returns a \texttt{NEGATIVE\_AMORTIZATION} warning. If the debt is
not cleared within the default simulation limit of 600 months, the plan
is infeasible within the horizon. The formula lets FR-10 separate the
computation clearly from the LLM's interpretation.

\section{Evaluation Methodology}

\subsection{Extraction and Classification Accuracy}

For a field to be extracted, let $TP$ be the correctly predicted values,
$FP$ the wrongly predicted or extraneous values, and $FN$ the
ground-truth values that were missed:

\begin{equation}
  Precision=\frac{TP}{TP+FP},\qquad
  Recall=\frac{TP}{TP+FN},\qquad
  F1=\frac{2\,Precision\,Recall}{Precision+Recall}.
\end{equation}

An exact match is achieved only when every required field of a
transaction is simultaneously correct. Macro-F1 averages the F1 of each
category so that a many-sample class does not mask a few-sample class
\cite{sokolova2009}. These metrics must be computed on an evaluation set
independent of the development prompts or rules.

\subsection{OCR/STT Error and Business Impact}

Let $S$, $D$, $I$ be the number of substitutions, deletions, and
insertions, and $N$ the number of characters or words in the ground
truth:

\begin{equation}
  CER=\frac{S_c+D_c+I_c}{N_c},\qquad
  WER=\frac{S_w+D_w+I_w}{N_w}.
\end{equation}

CER/WER measure raw text; transaction-field accuracy is what reveals
whether a media error corrupted the amount, date, merchant, or line
item. Therefore a provider smoke test must not be presented as accuracy.

\subsection{Algorithmic Correctness, Grounding, and Performance}

Deterministic algorithms are checked against a hand-computed expected
value, the absolute error $AE=|\hat y-y|$, and invariants such as the
conservation of the total across a wallet transfer. For narration, let
$N_{all}$ be the number of numeric expressions in the response and
$N_{supported}$ the number that can be mapped to a fact or an equivalent
formatting:

\begin{equation}
  Numeric\ Faithfulness=\frac{N_{supported}}{N_{all}}.
\end{equation}
\begin{equation}
  Unsupported\ Rate=1-Numeric\ Faithfulness.
\end{equation}

This is a \textbf{target evaluation method}; a checker covering every
number is not yet fully implemented. Latency uses percentiles: p50 is
the median and p95 is the value that 95\% of samples do not exceed; it
must separate text/media and cache hit/miss and must record the
provider, the hardware, and the number of runs.

\section{The Controlled, Supporting Role of the LLM}

\subsection{Function Calling Instead of Delegating Business Authority}

The Transformer uses attention to model relationships between tokens and
is the foundation of the modern LLM \cite{vaswani2017}; the Gemini family
extends this to multimodal input \cite{gemini2023}. PERFIN takes
advantage of the ability to understand diverse phrasings but does not use
the model to compute a ledger or a statistic. Tools are declared with a
schema \cite{geminifunctioncalling}; the LLM proposes a tool name and
its arguments, and the application validates, builds a preview, waits for
confirmation, and only then executes. Structured output guarantees only
the shape of the JSON, not that the category is right, that the two
wallets differ, or that the date is within the valid domain.

\subsection{Responsibility Boundary and Grounding}

\begin{table}[htbp]
  \centering
  \caption{Responsibility boundary between the deterministic core and the LLM}
  \label{tab:llm-responsibility}
  \small
  \begin{tabularx}{\textwidth}{|p{3.8cm}|p{4.0cm}|X|}
    \hline
    \textbf{Task} & \textbf{Responsible component} &
    \textbf{Reason} \\ \hline
    Total income/expense, balance, limit & SQL and business service &
    Accurate and reproducible \\ \hline
    Trend, anomaly, runway, correlation & Deterministic algorithm
    functions & Have a formula and an independent test \\ \hline
    Budget and goals & Budget/Goal Planner & Have constraints and a
    simulation \\ \hline
    Understanding dates, currency units, multiple intents & LLM or local
    parser & A language/context problem \\ \hline
    Selecting a tool and filling arguments & LLM, then the service
    checks & A bridge from an utterance to an API \\ \hline
    Writing interpretation/persona & LLM or template fallback & Must not
    change the facts \\ \hline
    Write, delete, transfer, re-tag & Service after confirmation & Have a
    data effect and need an audit \\ \hline
  \end{tabularx}
\end{table}

Insight facts must include the value, the unit, the observation window,
the number of samples, the method, and a missing-data warning. A persona
may change the tone in the spirit of behavioral design \cite{thaler2009}
but must not change the facts. The system currently separates the facts
from the narrator and has a template fallback; a comprehensive numeric
checker is still a target design that needs evaluation. The LLM does not
connect directly to PostgreSQL and does not execute a tool on its own.

\subsection{Conditions Under Which Using an LLM Is Meaningful}

An LLM is only valuable if it improves a measurable metric over a form
and a local parser: a higher field F1 or exact match, fewer
clarification turns or steps, while latency, cost, and the
unsupported-number rate stay within acceptable thresholds. An ablation
must use the same data and the same task. If it does not create a
benefit large enough to offset the risk, the direct flow or the local
parser is preferred. This approach keeps the focus of the project on
data and algorithms, with the LLM as a controlled communication layer.

\section{Infrastructure, Technology, and Selection Rationale}

\begin{table}[htbp]
  \centering
  \caption{Technology, selection criteria, and alternatives}
  \label{tab:technology-rationale}
  \small
  \begin{tabularx}{\textwidth}{|p{3.2cm}|X|X|}
    \hline
    \textbf{Technology} & \textbf{Why it fits the problem} &
    \textbf{Boundary or alternative} \\ \hline
    React Native + Expo & One codebase for text, image, and audio input
    and preview on mobile & Contains no formula; native code is needed
    only for device-specific requirements \\ \hline
    Node.js + Express & I/O-oriented API; shares JavaScript with the pure
    functions and their tests; fits a modular monolith & Avoids the
    operating cost of microservices \\ \hline
    PostgreSQL & Foreign keys, \texttt{DECIMAL}, transactions, and
    aggregate queries fit a ledger & A file or NoSQL store struggles to
    guarantee the same constraints and rollback \\ \hline
    Redis & TTL, cache-aside, pending state, and queue & Not a source of
    truth; the memory fallback is for development only \\ \hline
    BullMQ & Scheduling, retries, and a job ID for idempotent proactive
    tasks & A simple cron lacks the same queue/retry mechanism \\ \hline
    Gemini or local parser & The LLM widens language coverage; the parser
    guarantees a testable fallback & The provider sits behind an adapter
    and does not lock the business logic \\ \hline
    Local or cloud OCR/STT & Enables comparison of accuracy, latency,
    cost, and privacy & Raw output always passes the same
    validation/preview \\ \hline
  \end{tabularx}
\end{table}

Redis/BullMQ serve the recurring reminder, the runway scan, the
subscription scan, the month-end report, and the export cleanup. A job
must be small and idempotent; a unique event key in the internal message
prevents a duplicated effect on retry \cite{bullmq2026}. Scheduled
auto-backup does not yet have a corresponding worker and must not be
regarded as operational.

\section{Related Research and Applications}

Linear regression, z-score, IQR, and correlation are explainable
techniques that fit the goals of the project because their output can be
checked independently. OCR, STT, and multimodal LLMs widen the input
channels but do not by themselves guarantee the correctness of a
financial record.

Existing personal-finance systems tend to fall into three groups. The
first is the form/dashboard application. International products such as
Mint \cite{mint2023} and You Need A Budget (YNAB) \cite{ynab2024}, and
Vietnamese products such as Money Lover \cite{moneylover2024} and the
MISA money-management application, offer structured data and rich
reports but rely on multi-step manual entry or on bank synchronization
that is unavailable to most Vietnamese users. The second group is the
natural-language chatbot, which lowers entry friction but, when it lets
the model both extract and compute, produces numbers that are hard to
trace and audit. The third group is the specialized analytical model,
which is strong at detecting patterns but usually omits the
confirmation loop that a non-expert user needs before data is written.
On the extraction side, sequence-labeling approaches to entity
extraction \cite{lample2016} inform the field-level metrics used here,
but they require labeled data that PERFIN treats as future work rather
than as an achieved result.

PERFIN does not compete on commercial completeness. The gap it targets
is the combination of a relational ledger, deterministic algorithms, a
human in the loop, and a bounded language layer, evaluated with a
protocol that keeps the model out of the computation path.

\begin{table}[htbp]
  \centering
  \caption{Gaps addressed by PERFIN and how each is handled}
  \label{tab:research-gaps}
  \small
  \begin{tabularx}{\textwidth}{|X|X|X|}
    \hline
    \textbf{Gap} & \textbf{Risk} &
    \textbf{How it is handled} \\ \hline
    Natural input but easy to record wrongly & Dirty data spreads into
    reports & Typed draft, validation, preview, confirm \\ \hline
    Media has raw text but uncertain fields & Wrong amount, date, or
    category & Measure CER/WER and field accuracy; confirm before write
    \\ \hline
    Chatbot returns numbers of unclear origin & Hallucination and hard to
    audit & Query tools, JSON facts, a target grounding checker \\
    \hline
    Dashboard has numbers but lacks interpretation & Hard to notice a
    pattern & Deterministic algorithms; the narrator only reads facts \\
    \hline
    Classification does not adapt & Repeats the same error & Correction
    retrieval, a fuzzy threshold, and a margin \\ \hline
    Retry creates duplicate notifications & Bad experience and wrong data
    & Job ID, unique event key, idempotent handler \\ \hline
  \end{tabularx}
\end{table}
